\section{Conclusion}

We presented \textbf{Amphion EmoAnalyzer}, a compact 1.5B-parameter bilingual speech emotion understanding model that unifies categorical speech emotion recognition (SER) with descriptive emotion analysis (DEA). By combining multi-task speech-language training, a curated 1,000-hour bilingual emotion corpus, and DEA-focused Emotional DPO, EmoAnalyzer moves beyond closed-set emotion labels toward grounded, natural-language descriptions of how emotion is expressed in speech.

Across five open-source SER benchmarks, EmoAnalyzer achieves the strongest average performance in our unified evaluation setting, demonstrating robust categorical emotion recognition across English and Chinese datasets. On the DEA benchmark, EmoAnalyzer also outperforms leading proprietary MLLMs, including Gemini-3.0-Flash, showing that targeted emotion specialization can deliver stronger paralinguistic fidelity even with a compact model backbone.

EmoAnalyzer's design reflects a broader principle: emotion understanding in speech is fundamentally a multimodal alignment problem that requires explicit training connections among semantic grounding (ASR), categorical recognition (SER), and language generation (DEA). General-purpose audio-language models acquire these capabilities incidentally, if at all; targeted co-training and preference alignment make them robust. As speech models are increasingly deployed in real-world scenarios---customer service analytics, mental health monitoring, affective media analysis---this alignment will become a necessary property of practical audio-language systems rather than a specialized research objective.

Looking forward, we plan to further expand EmoAnalyzer's coverage across languages, speaking styles, and real-world acoustic environments. We are also interested in improving controllable emotion reasoning, strengthening long-context conversational emotion tracking, and releasing broader demos and evaluation resources to support practical emotion-aware speech applications.
